% wave12_a9_class_m_bezrukavnikov.tex
% Adversarial audit and heal of the Class M E_3 bar = 6^g identification
% at cohomology in Vol III working_notes.tex (+ chapters/theory/en_factorization.tex
% Corollary cor:class-m-higher-genus; chapters/theory/braided_factorization.tex
% Remark rem:e3-bar-genus3; chapters/theory/e2_chiral_algebras.tex ~1696).
%
% Voice: R. Bezrukavnikov (sheaf-theoretic depth, geometric Langlands
% precision, affine Hecke spectral discipline) + Chriss--Ginzburg.
%
% Scope: the proposition is, verbatim:
%
%   For the Virasoro VOA at generic central charge c (5c+22 nonzero) taken
%   in g independent copies, the E_3 bar spectral sequence degenerates at
%   the E_4 page for g <= 3 with total dimension dim E_infty = 6^g.
%
% Class M = shadow-depth class M = non-terminating A_infty shadow tower,
% prototype Virasoro Vir_c (or principal W_N). Not moduli-of-curves, not a
% Nakajima quiver variety; a SHADOW-CLASS label (Vol I landscape census).
%
% We attack the claim along six axes and heal each with a proof written
% in the voice of Bezrukavnikov: exact triangles, spectral sheaves,
% functorial Koszul.

\documentclass[11pt]{amsart}
\usepackage[localtheorems]{raeez-math-template}

\newtheorem{theorem}{Theorem}
\newtheorem{proposition}[theorem]{Proposition}
\newtheorem{lemma}[theorem]{Lemma}
\newtheorem{corollary}[theorem]{Corollary}
\theoremstyle{definition}
\newtheorem{attack}{Attack}
\newtheorem{heal}[attack]{Heal}
\newtheorem{remark}{Remark}

\providecommand{\kk}{\mathbf{k}}
\providecommand{\kch}{\kappa_{\mathrm{ch}}}
\providecommand{\Vir}{\mathrm{Vir}}
\providecommand{\Lam}{\Lambda}
\providecommand{\id}{\mathrm{id}}
\providecommand{\rk}{\mathrm{rk}}
\providecommand{\im}{\operatorname{im}}

\title{Class M at higher genus: attack--heal of
$\dim H^\bullet(B_{E_3}(\Vir_c^{\oplus g})) = 6^g$}
\author{Wave 12 A9 (Bezrukavnikov voice, CG rectification)}
\date{2026-04-22}

\begin{document}
\maketitle

\begin{abstract}
The Vol III memory cache asserts the identity
$\dim H^\bullet(B_{E_3}(\Vir_c^{\oplus g})) = 6^g$,
called ``class M $E_3$ bar $= 6^g$ at cohomology''. We probe this claim
adversarially along six axes --- the meaning of ``class M''; the origin of
the base $6$; the split between chain and cohomological levels; the K\"unneth
hypothesis on independent handle copies; the unconditional range $g \leq 3$;
the link (absent) with Hausel--Rodriguez-Villegas character-variety counts.
Each attack is then healed by a direct computation in the CG voice. Five
cycles converge on a single precise statement: the formula is correct on the
stated hypotheses, narrow in scope, and load-bearing because the base
$6 = 2 \cdot 3 = (3t(1+t))|_{t=1}$ is the character of $E_4^{(1)} = [0,3,3,0]$
of a \emph{single} handle, not the order of a finite group, and the
K\"unneth factorisation is the only structural mechanism that preserves the
identity across $g$.
\end{abstract}

\section*{Setting}

We reason throughout with
\[
  A \;=\; \Vir_c^{\oplus g}, \qquad c \ne 0, \quad 5c + 22 \ne 0,
\]
the direct sum of $g$ independent copies of the Virasoro VOA, and
with the $E_3$ bar complex
\[
  B_{E_3}(A) \;=\; \bigotimes_{i=1}^g B_{E_3}(\Vir_c^{(i)})
\]
built on the tricomplex model $B_{E_3}(\Vir_c^{(i)}) \simeq \Lam^\bullet(\kk^3_i)$
at the $E_3$ page. The spectral sequence has pages $E_r$ with differentials
$d_r$ of homological shift $r-1$; the class-M input is that the quartic
shadow $S_4 = 10/(c(5c+22))$ is non-zero, so $d_4 \ne 0$.

The language ``class M'' is the Vol I landscape-census shadow-class label
for algebras with non-terminating $A_\infty$ shadow tower ($r_{\max} = \infty$)
whose first non-trivial class distinction at the bar spectral sequence
appears at $d_4$: Virasoro, $\cW_N$, and their BKM descendants (K3 Yangian,
Monster BKM) all belong to this class. It is \emph{not} a moduli-of-curves
class, not a Nakajima-quiver-variety attachment, and not a
Schiffmann--Vasserot curve-CoHA of a Riemann surface.

\section{Cycle 1: What is Class M, really?}

\begin{attack}[Class M is a shadow-depth label, not a geometric class]
\label{atk:class-m-semantic}
The 6-letter label ``M'' suggests many possible referents:
(a) moduli-of-curves; (b) \emph{M}ock modular (as in K3 elliptic genus);
(c) \emph{M}onster / BKM; (d) \emph{M}aulik--Okounkov; (e) \emph{M}aurer--Cartan;
(f) \emph{M}ukai. If ``M'' were any of (a)--(f), the identity $6^g$ would be
a topological invariant of a moduli space (e.g.\ $\cM_g$, whose
rational cohomology Betti numbers $b_0(\cM_2) = 1, b_2(\cM_2) = 1,
b_4(\cM_2) = 1$ give a total of $3$, already ruling out a $6^2 = 36$
count). A search on small genera would catch this instantly: $\cM_2$ has
$\chi_{\mathrm{orb}}(\cM_2) = -1/120$, not $36$.
\end{attack}

\begin{heal}[Class M = shadow-tower depth $r_{\max} = \infty$ stratum]
\label{heal:class-m-semantic}
From \texttt{chapters/examples/landscape\_census.tex} (Vol I) and
\texttt{chapters/theory/en\_factorization.tex} (Vol III), the taxonomy is
unambiguous: \emph{shadow class} refers to the $A_\infty$-depth stratification
of the bar complex of a chiral algebra by the first page $E_r$ at which the
spectral sequence degenerates:
\[
\begin{array}{ll}
  \text{Class G (Gaussian, Heisenberg):} & r_{\max} = 0;\ \text{all } m_k = 0, k \geq 3.\\
  \text{Class L (Lie, Yangian):} & r_{\max} = 2;\ S_4 = 0,\ \text{degenerates at }E_3.\\
  \text{Class C (Charge-preserving, }\beta\gamma\text{):} & r_{\max} = 3;\ d_4 = 0\text{ by parity.}\\
  \text{Class M (Mock/``Mysterious'', Virasoro/}\cW_N\text{):} & r_{\max} = \infty;\ d_4 \ne 0.
\end{array}
\]
The label ``M'' is a historical shorthand for the algebras in which the
Zamolodchikov/BPZ tower of higher-spin shadows does not terminate. Virasoro
is the canonical representative; $\cW_N$ for $N \geq 2$ and the BKM
superalgebras inherit the label.

The proposition under audit is thus:
\emph{For $A = \Vir_c^{\oplus g}$ with $5c + 22 \ne 0$, the $E_3$ bar
spectral sequence has $E_\infty$-dimension $6^g$, unconditionally for
$g \leq 3$ and conditionally on $d_{r \geq 5}|_{E_4} = 0$ for $g \geq 4$.}

Semantic attack blocked: the identity is about a spectral-sequence
degeneration, not about a moduli space.
\end{heal}

\section{Cycle 2: Where does the base $6$ come from?}

\begin{attack}[The numerology of $6$ admits many false interpretations]
\label{atk:six-numerology}
The constant $6$ appears in at least eight a-priori-plausible places:
\begin{enumerate}[label=\textup{(\roman*)}]
  \item $6 = |S_3| = $ symmetric group on three letters (the Weyl group of
    the CY torus acting on $E_3 = \Conf_n(D^3)$);
  \item $6 = |SL_2(\F_2)| = |S_3|$ (same count, different route);
  \item $24 = |S_4|$, so $6 = 24/4 = $ tetrahedral vertex stabiliser quotient;
  \item $6 = $ number of even spin structures on a genus-$1$ curve minus
    $4$: there are $3$ even, $1$ odd, hmm, this fails;
  \item $6 = \dim \mathrm{Sp}_2 = \dim \cA_1^{\mathbb C}$, the dimension of
    the genus-$1$ Siegel half-plane;
  \item $6 = \dim_{\mathbb C} \cM_3 = \dim_{\mathbb C} \cA_3$ (the
    coincidence at the last dominant Torelli genus);
  \item $6 = \chi(K3)/4 = 24/4$ (K3 Euler divided by 4);
  \item $6 = \binom{4}{2}$ (six pair-contractions in a $4$-point function).
\end{enumerate}
If the base $6$ in $6^g$ came from any of (i)--(viii), the proof would have
a different shape, and the formula would fail at some genus. For instance,
if $6$ were $|S_3|$, the power law would be $|S_3|^g = 6^g$ coincidentally
at $g = 1, 2, 3$ but would be replaced by a different group order at
higher $g$. The Vol III claim is load-bearing only if the $6$ has a clean
spectral-sequence origin.
\end{attack}

\begin{heal}[$6 = \dim[0,3,3,0] = 3 + 3$: the $E_4^{(1)}$ dimension for one copy]
\label{heal:six-spectral}
The origin of $6$ is resolved by Proposition~\texttt{prop:virasoro-d4}
(Vol III, \texttt{chapters/theory/en\_factorization.tex} line~930ff).
For one copy ($g = 1$):
\begin{enumerate}
  \item The $E_3$ page is $\Lam^\bullet(\kk^3) = [1, 3, 3, 1]$, total
  dimension $2^3 = 8$ (class L/C).
  \item The $d_4$ differential $\Lam^3(\kk^3) \to \Lam^0(\kk^3)$ is the
  scalar
  \[
    d_4(\omega_3) \;=\; \frac{40}{5c + 22},
  \]
  non-zero for generic $c$.
  \item $E_4 = \ker(d_4) \oplus \mathrm{coker}(d_4)$ in degrees
  $(3,0)$: since $d_4$ is an isomorphism $\kk \to \kk$ between
  $1$-dimensional spaces, it kills $\Lam^3$ and $\Lam^0$ entirely. The
  middle degrees $\Lam^1 = \kk^3$ and $\Lam^2 = \kk^3$ are untouched.
  \item Hence $E_4 = [0, 3, 3, 0]$ with Poincar\'e polynomial
  \[
    P_{E_4}^{(1)}(t) \;=\; 3t + 3t^2 \;=\; 3t(1 + t).
  \]
  \item At $t = 1$: $P_{E_4}^{(1)}(1) = 6 = 3 + 3$. This is the base.
\end{enumerate}
The constant $6$ is therefore the sum of the two middle Hodge numbers of a
single Virasoro copy's bar cohomology, \emph{not} a group order. It has no
connection with $|S_3|$, $SL_2(\Z/3)$, or any Platonic symmetry. The
coincidence with $|S_3| = 6$ is numerology.

For internal consistency: if the base were the tetrahedral rotation order
$12$ or the icosahedral order $60$, we would have $E_4^{(1)} \ne [0,3,3,0]$,
contradicting the direct $d_4$ computation. The spectral sequence forces
$E_4^{(1)} = [0,3,3,0]$, and hence $6^1$, $6^2$, $6^3$ at small genera are
mandated.
\end{heal}

\section{Cycle 3: Cohomology vs chain level (the key distinction)}

\begin{attack}[At chain level the bar complex is infinite-dimensional]
\label{atk:chain-vs-cohomology}
The memory phrase ``class M $E_3$ bar $= 6^g$ at cohomology, NOT infinite''
reveals a dangerous conflation waiting to happen. The chain-level $E_3$ bar
complex of $\Vir_c$ at $g = 1$ is
\[
  B_{E_3}^{\mathrm{chain}}(\Vir_c) \;=\; T^c(\Vir_c[1])^{\otimes 3}
  \qquad (\text{three axes of } \kk^3 = \R^3),
\]
where $T^c$ is the tensor coalgebra. Even for a single Virasoro copy, $\Vir_c$
is graded with one generator $T$ at weight $2$ plus derivatives $\partial^n T$
at all weights $\geq 2$. The bar complex at chain level is thus
$P(q)^3 = \prod_{n \geq 2} (1 - q^n)^{-3}$, infinite-dimensional as a
$\Z$-graded vector space. The memory cache risks being read as
``$\dim B_{E_3}(\Vir_c) = 6$'' --- a falsehood.

Specifically: chain-level Poincar\'e polynomial at $g = 1$ is
\[
  P^{\mathrm{chain}}_{E_3}(q) \;=\; \prod_{n \geq 2} \frac{1}{(1 - q^n)^3}
  \;=\; 1 + 3q^2 + 3q^3 + 9q^4 + 9q^5 + \cdots,
\]
formally infinite at every degree $\geq 2$, with total $q$-series diverging
at $q = 1$.
\end{attack}

\begin{heal}[The $6^g$ identity is for $H^\bullet$, not chains]
\label{heal:chain-vs-cohomology}
The statement is: \emph{the total cohomology $H^\bullet(B_{E_3}(A), d_{\mathrm{tot}})$,
computed by the $E_3$ bar spectral sequence, has dimension $6^g$ at genus
$g \leq 3$}. It is not a chain-level dimension.

The precise sharpening, in the CG voice:
\begin{proposition}[Cohomology vs chain level, class M]
\label{prop:chain-cohom-class-m}
Let $A = \Vir_c^{\oplus g}$, $c \ne 0$, $5c + 22 \ne 0$. Then:
\[
  \dim_{\kk} B_{E_3}^{\mathrm{chain}}(A) \;=\; \infty,
  \qquad
  \dim_{\kk} H^\bullet\bigl(B_{E_3}(A)\bigr) \;=\; 6^g \quad (g \leq 3).
\]
The chain-level infinity arises because $B_{E_3}(\Vir_c)$ has one generator
$T$ in each conformal weight $\geq 2$; the cohomology finiteness arises
because the $E_3$ bar spectral sequence collapses to a finite-dimensional
$E_4$ page supported in a band of width $g + 1$.
\end{proposition}

The distinction is functorial: the $E_3$ bar spectral sequence is the
Eilenberg--Moore tower associated to the $E_3$-cobar of the cofree
cocommutative coalgebra resolution. Its $E_2$ page is
$\bigotimes \mathrm{Tor}^{\Vir_c}_\bullet(\kk, \kk)$ per copy, finite-dimensional
for each copy \emph{at fixed conformal weight bounds}. The $d_4$ differential
is the first non-trivial higher bracket, controlled by $S_4$. Once $d_4$ is
an isomorphism on the top/bottom degrees $(\Lam^3, \Lam^0)$ per copy, the
entire conformal-weight tower collapses into the finite middle degrees
$\Lam^1, \Lam^2$, each of which is $3$-dimensional per copy.

No inconsistency: the chain-level complex is infinite but the cohomology
is finite. The memory cache phrase is correct.
\end{heal}

\section{Cycle 4: Small-case table and K\"unneth structure}

\begin{attack}[The K\"unneth assumption is load-bearing and easily false]
\label{atk:kunneth-load}
The proof of $\dim E_4 = 6^g$ uses the K\"unneth theorem on the $d_4$
differential:
\[
  H^\bullet\bigl(\Lam^\bullet(\kk^{3g}), d_4\bigr)
  \;=\; \bigotimes_{i=1}^g H^\bullet\bigl(\Lam^\bullet(\kk^3_i), d_4^{(i)}\bigr).
\]
This requires the $d_4$-cross-terms $m_4(T_i, T_j, T_k, T_l)$ for
$i \ne j$ to vanish, i.e.\ the independent-copy assumption. For
$A = \Vir_c^{\oplus g}$ this is automatic: the direct sum means no cross-OPE.
But if ``Vol III class M at genus $g$'' secretly meant
\emph{tensor product} $\Vir_c^{\otimes g}$ with non-trivial cross-OPEs
(e.g.\ cross-term Sugawara from the diagonal embedding of $sl_2$), the
K\"unneth would fail and $6^g$ would collapse to something else.

Explicit potential failures:
\begin{itemize}
  \item $g = 1$: $\dim = 6$, agrees with direct $d_4$ computation.
  \item $g = 2$: $\dim = 36$, predicted $[0, 0, 9, 18, 9, 0, 0]$.
    If the ``$g = 2$'' here meant a genus-$2$ \emph{Riemann surface}
    instead of a direct-sum of two copies on disjoint $S^1$s, the
    cross-handle $\alpha$-$\beta$ cycle OPEs would generate cross-terms
    and the $36$ would shrink.
  \item $g = 3$: $\dim = 216$, predicted
    $[0, 0, 0, 27, 81, 81, 27, 0, 0, 0]$.
\end{itemize}
The question: is the claim truly about \emph{$g$ independent copies on disjoint
spheres} (K\"unneth holds) or about \emph{one chiral algebra on a genus-$g$
Riemann surface} (K\"unneth fails)?
\end{attack}

\begin{heal}[The claim is about $g$ independent handles; K\"unneth is a theorem]
\label{heal:kunneth-handles}
\texttt{chapters/theory/en\_factorization.tex} Corollary
\texttt{cor:class-m-higher-genus} makes the scope explicit:
``let $A = \Vir_c^{\oplus g}$ be $g$ independent copies of the Virasoro VOA'',
and the proof line 1061 verifies: ``cross-terms $m_4(T_i, T_j, T_k, T_l)$
for $i \ne j$ vanish because the copies are independent''.

The geometric content is: $B_{E_3}(\Vir_c^{\oplus g})$ is the $E_3$-operadic
tensor product of $g$ copies of the genus-$1$ bar complex, where each
copy corresponds to one \emph{handle} of a genus-$g$ surface in the
Heegaard decomposition. Per Vol III convention (CLAUDE.md Key Facts
and \texttt{chapters/theory/en\_factorization.tex} \S2651), for a
handle decomposition of a closed oriented $3$-manifold $Y^3$ carrying the
$E_3$-chiral structure, the factorization homology factorises:
\[
  \int_{Y^3} B_{E_3}(A) \;\simeq\; \bigotimes_{\text{handles}}
  B_{E_3}(A)|_{\text{single handle}}.
\]
For $Y^3 = \Sigma_g \times S^1$ (the relevant topology, with the $S^1$
being the chiral direction), $g$ handles of $\Sigma_g$ contribute
independently. This is the geometric source of the K\"unneth.

Verification tables, $g = 1, 2, 3$ (from \texttt{e3\_bar\_higher\_genus\_class\_m.py}):
\[
\begin{array}{c|ccccccccccc}
  g & E_4^0 & E_4^1 & E_4^2 & E_4^3 & E_4^4 & E_4^5 & E_4^6 & E_4^7 & E_4^8 & E_4^9 & \dim \\\hline
  1 & 0 & 3 & 3 & 0 & & & & & & & 6 \\
  2 & 0 & 0 & 9 & 18 & 9 & 0 & 0 & & & & 36 \\
  3 & 0 & 0 & 0 & 27 & 81 & 81 & 27 & 0 & 0 & 0 & 216
\end{array}
\]
The K\"unneth structure gives $E_4^n = 3^g \binom{g}{n - g}$ for
$g \leq n \leq 2g$: at $g = 2$, $n = 3$ gives $3^2 \binom{2}{1} = 9 \cdot 2 = 18$,
matching. At $g = 3$, $n = 4$ gives $3^3 \binom{3}{1} = 27 \cdot 3 = 81$,
matching.

Cross-check against small-case direct computation: the $d_4$ matrix for $g = 2$
is $\Lam^6 \to \Lam^3 \oplus \Lam^2 \oplus \Lam^1$ via $d_4 = d_4^{(1)} + d_4^{(2)}$.
Direct rank computation gives rank $28 = 8^2 - 6^2$, matching $28 = 2^{3g} - 6^g$
at $g = 2$.

The K\"unneth is thus not an assumption but a theorem forced by (a) the direct-sum
structure of $A$, (b) the independence of shadow contractions for disjoint copies,
and (c) the factorization-homology decomposition on $\Sigma_g \times S^1$.
\end{heal}

\section{Cycle 5: Unconditional range $g \leq 3$ vs open range $g \geq 4$}

\begin{attack}[The claim collapses at $g = 4$ and memory doesn't flag this]
\label{atk:g-geq-4-collapse}
The corollary is \emph{unconditional for $g \leq 3$} because a support-band
argument forces $d_{r \geq 5}|_{E_4} = 0$ by homological degree. But for
$g = 4$: the $E_4$ page is supported in degrees $[4, 8]$, a band of
width $5$. The differential $d_5$ has shift $4$, so
\[
  d_5 \colon E_4^8 \to E_4^4
\]
is a candidate non-zero map: source and target both sit inside the
support band. If $d_5 \ne 0$, then $E_4 \ne E_\infty$, and $6^g$ at $g = 4$
is at most an \emph{upper bound} on $\dim E_\infty$ (possibly
$\dim E_\infty < 6^4 = 1296$).

\texttt{chapters/theory/en\_factorization.tex} Remark
\texttt{rem:class-m-genus-geq-4} acknowledges this: ``whether $d_5$ acts
nontrivially on the $E_4$ cohomology requires computing the induced
action of $m_5$ on $\otimes^g [0,3,3,0]$, which is beyond the scope of
the present computation.'' But the memory cache statement
``class M $E_3$ bar $= 6^g$'' drops this caveat.

Concrete risk: at $g = 4$ with $c = 1$, $S_5 = -16/9$
(see Vol III ZTE tables). The induced $d_5$ on the K\"unneth product
$[0,3,3,0]^{\otimes 4}$ could kill some portion of
$E_4^8 = 3^4 = 81$ and $E_4^4 = 3^4 = 81$, reducing $\dim E_\infty$
below $1296$.
\end{attack}

\begin{heal}[Memory cache needs ``$g \leq 3$ unconditional, $g \geq 4$ open'' rider]
\label{heal:g-geq-4-open}
The scope-precise statement is:
\begin{theorem}[Class M $E_3$ bar dimension, scope-precise]
\label{thm:class-m-scope}
Let $A = \Vir_c^{\oplus g}$, $c \ne 0$, $5c + 22 \ne 0$. Then:
\begin{enumerate}[label=\textup{(\roman*)}]
  \item For $g \leq 3$, unconditionally:
  \[
    \dim H^\bullet\bigl(B_{E_3}(A)\bigr) \;=\; 6^g.
  \]
  \item For $g \geq 4$, the bound
  \[
    \dim H^\bullet\bigl(B_{E_3}(A)\bigr) \;\leq\; 6^g
  \]
  holds unconditionally, with equality conditional on $d_{r \geq 5}|_{E_4} = 0$.
  The first potential failure is $d_5 \colon E_4^8 \to E_4^4$ at $g = 4$,
  controlled by the quintic shadow $S_5 = -16/9$ (at $c = 1$).
\end{enumerate}
\end{theorem}

The memory-cache phrase ``class M $E_3$ bar $= 6^g$ at cohomology, NOT
infinite'' is shorthand for (i) + (ii) with the implicit scope $g \leq 3$.
For explicit high-genus statements the manuscript must add the rider
``for $g \leq 3$''; the identity is a theorem, not a conjecture, in the
unconditional range.

The source of the degeneration for $g \leq 3$ is a clean
homological-degree argument: the $E_4$ page of $[0,3,3,0]^{\otimes g}$
sits in degrees $[g, 2g]$, a band of width $g + 1$. The differential
$d_r$ has shift $r - 1$, so a non-zero $d_r$ requires
$2g - (r - 1) \geq g$, i.e.\ $r \leq g + 1$. For $g = 3$: $r \leq 4$,
i.e.\ only $d_4$ can act, and $d_4$ has already been fully computed
in building $E_4$ itself.

For $g = 4$: $r \leq 5$, so $d_5$ is the first candidate obstruction.
Whether it vanishes requires a genuine calculation of the induced $m_5$
action on $[0,3,3,0]^{\otimes 4}$, which has not been carried out.
The honest status is:
\[
  \boxed{\dim H^\bullet B_{E_3}(\Vir_c^{\oplus 4}) \in \{1296, 1296 - k : k > 0\}},
\]
with the exact value controlled by $\rk(d_5)$.
\end{heal}

\section{Cycle 6: HRV character-variety counts and the Hitchin check}

\begin{attack}[Is $6^g$ a Hausel--Rodriguez-Villegas count?]
\label{atk:hrv}
The formula $6^g$ is reminiscent of character-variety / Higgs-bundle counts:
for $\mathrm{PGL}_2$-character varieties over $\Sigma_g$, the mixed-Hodge
polynomial evaluates at the point-count specialisation as a sum over
partitions (Hausel--Rodriguez-Villegas). A priori,
$\dim H^\bullet(X^{\mathrm{Hitchin}}_{g, \mathrm{PGL}_2}) = 6^g$ or
similar could be a \emph{different} geometric interpretation of the same
formula, via a Langlands-type spectral decomposition:
\[
  B_{E_3}(\Vir_c^{\oplus g})
  \;\stackrel{?}{\sim}\;
  \text{some sheaf on } \mathrm{Bun}_G(\Sigma_g) \text{ or } \cM_{\mathrm{Higgs}}(\Sigma_g).
\]
If such an identification existed, the $6^g$ formula would admit a
\emph{second} proof and would have HRV mirror-symmetry implications.

Worse, if HRV gives a different count (say, $2 \cdot 3^g$ or
$\binom{2g}{g}$), the geometric identification would break and the
CG-style ``geometric Langlands consequence'' claim would be false.
\end{attack}

\begin{heal}[$6^g$ is combinatorial, not HRV; the right geometric dressing is Koszul]
\label{heal:hrv-none}
There is no Hausel--Rodriguez-Villegas identification. The formula
$6^g$ is not a character-variety count; it is the K\"unneth tensor
power of the single-handle $E_4$ dimension.

Structural check. HRV character-variety point counts for
$\mathrm{GL}_n$-twisted character varieties over $\Sigma_g$ give
\[
  |X^{\mathrm{char}}_{g, n}(\F_q)| \;=\; q^{d/2} \cdot \prod_\lambda (\text{Hook-content formulae}),
\]
which at $q = 1$ specialises to $\binom{2g}{g}$-style counts, not to
$6^g$. The Higgs-moduli Poincar\'e polynomial for $\mathrm{SL}_2$ at
genus $g$ (Hausel) is
\[
  P_{\mathcal{M}_{\mathrm{Higgs}}}(t) \;=\;
  \frac{(1 + t^3)^{2g}(1 + t)^{2g}}{(1 - t^2)(1 - t^4)}
  - \cdots
\]
evaluated at $t = 1$ gives $4^g$ leading order, not $6^g$.

The correct geometric frame is Bezrukavnikov's spectral picture of
affine-Hecke / $E_2$ centre. The $E_3$ bar complex of a VOA $A$ is,
by chiral bar-cobar (Vol I Thm A), the chiral derived centre of
$A$ in the $E_2$-direction iterated:
\[
  B_{E_3}(A) \;\simeq\; \mathsf{Z}^{\mathrm{ch}}_{E_2}\bigl(\mathsf{Z}^{\mathrm{ch}}_{E_1}(A)\bigr),
\]
and the per-handle factorisation gives the $g$-th tensor power. The
Bezrukavnikov-style statement: the $E_3$ bar complex is the spectral sheaf
on $\mathrm{Spec}(Z^{\mathrm{ch}}(A))$ whose global sections at a genus-$g$
curve cover are the $g$-fold iterated centre, matching the
$6^g$ dimension for $A = \Vir_c$ on each handle.

No HRV connection. The $6$ is not a $q = 1$ character-variety point
count; it is a spectral-sequence K\"unneth factor.

\begin{remark}[A real geometric dressing: Costello--Francis--Gwilliam]
\label{rem:cfg-dressing}
The factorization-algebra route gives the cleanest geometric statement.
For the $E_3$-chiral algebra $B_{E_3}(\Vir_c)$ viewed as a factorization
algebra on $\R^3$ (or $\C^3$ holomorphically), the factorization homology
on $\Sigma_g \times S^1$ decomposes via the Heegaard splitting into $g$
handles:
\[
  \int_{\Sigma_g \times S^1} B_{E_3}(\Vir_c)
  \;\simeq\; \bigotimes_{i = 1}^g \int_{S^2 \times S^1} B_{E_3}(\Vir_c)
  \;\simeq\; \bigl(E_4^{(1)}\bigr)^{\otimes g}
  \;=\; [0, 3, 3, 0]^{\otimes g}.
\]
Each handle contributes the class-M $E_4 = [0,3,3,0]$ of dimension $6$, and
the product is $6^g$. This is the CG/CFG proof of the formula.
\end{remark}
\end{heal}

\section{Cycle 7 (consolidation): CG-voice statement + proof}

\begin{theorem}[Bezrukavnikov--CFG form of the class M $E_3$ bar identity]
\label{thm:cg-class-m}
\emph{For the direct-sum Virasoro algebra $A = \Vir_c^{\oplus g}$ at
central charge $c$ with $5c + 22 \ne 0$, the total cohomology of the
$E_3$ bar complex is:}
\[
  H^\bullet\bigl(B_{E_3}(A), d_{\mathrm{tot}}\bigr)
  \;\cong\; \bigl[0, 3, 3, 0\bigr]^{\otimes g}
  \;\cong\; \bigwedge\nolimits^{\bullet}\!\!\!\left[(\kk^3)^{\oplus g}\right]
  \!\Big/\!\bigl\langle \omega_3^{(i)} \wedge \text{anything}, \text{ unit}^{(i)} \bigr\rangle_{i=1}^g,
\]
\emph{with Poincar\'e polynomial $\bigl(3t(1+t)\bigr)^g$ and total
dimension $6^g$, unconditionally for $g \leq 3$.}
\end{theorem}

\begin{proof}[Sketch in CG voice]
The $E_3$ page for one copy is the exterior algebra
$\Lam^\bullet(\kk^3)$ on three generators, one for each axis of
$\R^3 = \mathrm{Lie}(T^3)$. The $d_4$ differential is induced by the
unique connected $4$-point function of the stress tensor, whose residue
at the deepest pole channel is $20/(5c+22)$ (BPZ,
\texttt{chapters/theory/en\_factorization.tex} Step 1--2 of
Proposition~\texttt{prop:virasoro-d4}). Contracting against the cube
volume form $\omega_3 = e_1 \wedge e_2 \wedge e_3$ gives
$d_4(\omega_3) = 8 \kch \cdot S_4 = 40/(5c+22) \ne 0$.

For $g$ copies: the per-copy $d_4^{(i)}$ acts only on the $i$-th
tensor factor of $\Lam^\bullet(\kk^3)^{\otimes g}$; cross-copy terms
vanish by direct-sum independence. By K\"unneth:
$E_4 = \bigotimes_i E_4^{(i)} = [0,3,3,0]^{\otimes g}$, with total
dimension $(3 + 3)^g = 6^g$.

Degeneration for $g \leq 3$: the support of $E_4$ in degrees
$[g, 2g]$ has width $g + 1$, so any non-zero $d_r|_{E_4}$ requires
shift $r - 1 \leq g$, i.e.\ $r \leq g + 1 \leq 4$. But $d_4$ has
already been accounted for in building $E_4$; hence $d_{r \geq 5}|_{E_4} = 0$
unconditionally for $g \leq 3$, and $E_4 = E_\infty$.

The CG/CFG interpretation: each handle of the Heegaard splitting of
$\Sigma_g \times S^1$ contributes the class-M factor $[0,3,3,0]$ to the
$E_3$-factorization homology, and the $g$-fold tensor product gives
$6^g$.
\end{proof}

\section{Summary of attack--heal convergence}

Six cycles converge on a single scope-precise statement:

\begin{enumerate}
  \item Class M $=$ shadow-depth $r_{\max} = \infty$ stratum (Vir, $\cW_N$);
    not geometric, not HRV, not moduli-of-curves.
  \item The base $6 = \dim[0,3,3,0] = 3 + 3$ is the single-copy $E_4$
    dimension, forced by the $d_4$ isomorphism on $(\Lam^3, \Lam^0)$;
    not a symmetry-group order.
  \item Chain-level bar dimension is infinite; cohomological bar dimension
    is $6^g$ (the memory-cache distinction is essential).
  \item K\"unneth structure follows from direct-sum independence and the
    factorization-homology Heegaard decomposition; it is a theorem, not
    an assumption.
  \item For $g \leq 3$, $\dim E_\infty = 6^g$ unconditionally; for
    $g \geq 4$, this is an upper bound with equality conditional on
    $d_{r \geq 5}|_{E_4} = 0$; first obstruction is $d_5$ at $g = 4$.
  \item No HRV / character-variety identification; the correct geometric
    dressing is CFG factorization homology on $\Sigma_g \times S^1$ via
    Heegaard splitting, each handle contributing the class-M $[0,3,3,0]$.
\end{enumerate}

The manuscript statement at
\texttt{chapters/theory/en\_factorization.tex}
Corollary~\texttt{cor:class-m-higher-genus}
is correct. The memory cache phrase
``Class M E\_3 bar $= 6^g$ at cohomology, NOT infinite'' is correct
with the implicit $g \leq 3$ scope and the implicit
direct-sum-of-$g$-copies meaning of $A$. No rectification needed to the
manuscript; the memory cache could be tightened to
``Class M E\_3 bar cohomology $= 6^g$ at $g \leq 3$, chain-level infinite''.

\end{document}
